\ProvidesFile{CascadeCall.tex}[Каскадный вызов]


\subsection{Каскадный вызов функций}

\subsubsection{Что мы хотим}

Если до или во время прочтения описания задачи вас преследуют мысли <<а нахрена оно мне надо?>>, то для самых любопытных сразу скажу, что нам придется решать эти задачи для автоматизации построения State Machine, которая будет обсуждаться в разделе~\ref{section::StateMachine}.

\paragraph{Описание задачи}

Предположим, что у нас есть некоторый фиксированный набор типов: \verb"A", \verb"B", \verb"C", \verb"D".
И есть функции, принимающие аргументы указанных типов (возможно по ссылкам) и возвращающие \verb"bool".
\begin{cppcode}[\nonumbers]
bool f0(A, B, C&, D);
bool f1(A, B, C, D);
bool f2(A&, B, const C&, D&);
\end{cppcode}
Теперь мы хотим вызвать эти функции на наборе переменных
\begin{cppcode}[\nonumbers]
A a; B b; C c; D d;
\end{cppcode}
следующим образом:
\begin{enumerate}
\item В начале вызываем
\begin{cppcode}[\nonumbers]
bool f0(a, b, c, d);
\end{cppcode}
Если она возвращает \verb"true", то останавливаемся и возвращаем \verb"0".
Иначе идем дальше.

\item Теперь вызываем
\begin{cppcode}[\nonumbers]
bool f1(a, b, c, d);
\end{cppcode}
Если она возвращает \verb"true", то останавливаемся и возвращаем \verb"1".
Иначе идем дальше.

\item Теперь вызываем
\begin{cppcode}[\nonumbers]
bool f2(a, b, c, d);
\end{cppcode}
Если она возвращает \verb"true", то останавливаемся и возвращаем \verb"2".
Иначе возвращаем \verb"-1".
\end{enumerate}
То есть мы хотим вызывать функции друг за другом на одном и том же наборе переменных до тех пор, пока функция не вернет \verb"true".
Если вернула, то мы останавливаемся и возвращаем номер функции, на которой остановились.
А если ни одна функция не вернула \verb"true", то возвращаем \verb"-1" как флаг того, что ничего не было найдено.

\paragraph{Какой хотим синтаксис}

В примере выше, мы бы хотели иметь возможность написать следующий код:
\begin{cppcode}[\nonumbers]
A a; B b; C c; D d;
using FList = AL::List<f0, f1, f2>;
int res = cascade_call<FList>(a, b, c, d);
\end{cppcode}
То есть мы складываем все функции в список элементов разного типа.
И потом вызываем функцию \verb"cascade_call", которая проходится по списку функций и каждую вызывает от своего набора аргументов.

\subsubsection{Идея имплементации \texttt{cascade\_call}}

\paragraph{Идея рекурсивного решения}

Понятно, как можно написать решение в рекурсивной форме.
Для этого нужно представить себе наш код как:
\begin{cppcode}
if (f0(a, b, c, d))
  return 0;
if (f1(a, b, c, d))
  return 1;
if (f2(a, b, c, d))
  return 2;
return -1;
\end{cppcode}
Тогда можно посмотреть на голову списка и вызывать ее. Если получили \verb"true", то возвращаем текущий индекс.
А если получили \verb"false", то увеличиваем текущий индекс на \verb"1", выкидываем голову списка и переходим к следующей функции.
В случае пустого списка выкинуть \verb"-1".
Я не буду писать это решение в коде%
\footnote{Если хотите, можете взять себе это как упражнение}
так как мне хочется иметь не рекурсивную версию.

\paragraph{Идея не рекурсивного решения}

Давайте вспомним, что в C++ логические операторы для \verb"bool" обладают ленивым вычислением.
Например выражение
\begin{cppcode}[\nonumbers]
f0(a, b, c, d) || f1(a, b, c, d) || f2(a, b, c, d);
\end{cppcode}
вычисляется слева направо до первого момента, как будет получено \verb"true".
После чего дальнейшие аргументы не вычисляются.
Это позволяет выполнить функции одну за другой и остановиться вовремя.
Минус в том, что нельзя так быстро понять, а на каком индексе мы остановились.
Для решения этой проблемы, мы заведем переменную \verb"ind" для индекса с изначальным значением \verb"-1", завернем наши функции во wrapper-ы которые будут принимать \verb"ind" по ссылке.
И каждая функция будет писать в \verb"ind" свой номер, если результат вычисления \verb"true".
А вычислить логическое <<или>> на всех функциях можно с помощью fold-expression доступных с C++17.
Давайте напишу условный код (так оно не сработает, конечно):
\begin{cppcode}[\nonumbers]
template<class... Funs, class... Args>
bool call_left_to_right(Funs... fs, Args&&... args) {
  return (... || fs(std::forward<Args>(args)...)); 
}
\end{cppcode}

\subsubsection{Имплементация \texttt{cascade\_call}}

Давайте я начну с кода самой функции.
Так как пока у меня нет рабочего решения, то идея заключается в том, чтобы внутри вызвать вспомогательный класс.
Причина почему класс, а не функция в том, что классы проще специализировать для проверки на ошибки.
\begin{cppcode}
namespace detail {
  ...
} // namespace detail

template<class FList, class... Args>
constexpr int cascade_call(Args&&... args) {
  return detail::CascadeImpl<FList>::call(std::forward<Args>(args)...);
}
\end{cppcode}
И вспомогательный шаблонный класс мы поместим во вспомогательное пространство имен, чтобы не засорять пространство имен.
Имплементация вспомогательного класса будет как обычно в два шага:
\begin{enumerate}
\item Общая версия с дефолтным параметром, которая отлавливает ошибки.

\item И специализированная версия, которая выполняет работу.
\end{enumerate}
Вот так выглядит общая версия:
\begin{cppcode}
template<class FList, bool isCorrect = AL::isList<FList>>
struct CascadeImpl {
  static_assert(
      AL::isList<FList>,
      "The template argument of cascade_call MUST be an element list of "
      "functions returning bool!");
};
\end{cppcode}
Тут я проверяю только, что \verb"FList" -- это список элементов вида \verb"List<u, v, w>", но не проверяю можно ли позвать \verb"u", \verb"v", \verb"w" с нужными аргументами.
Можно накручивать дополнительные проверки, как вам вздумается.
Тем более, что мы точно знаем как выглядят аргументы функций (с точностью до наличия ссылок или \verb"const" калификаторов).
Я этого делать здесь не буду.
Теперь общая имплементация выглядит так:
\begin{cppcode}
template<class FList>
struct CascadeImpl<FList, true> {
  template<class... Args>
  constexpr static int call(Args&&... args) {
    return CascadeUnpack<AL::Enum<FList>>::exec(std::forward<Args>(args)...);
  }
};
\end{cppcode}
Давайте поясню, что тут происходит.
\begin{enumerate}
\item Прежде чем позвать каждую функцию, мне надо знать ее номер.
Для этого я  список функций превращаю в список функций с индексами.
Для этого использую метод \verb"Enum", который обсуждался в разделе~\ref{section::TypeListOps}.%
\footnote{Формально мы обсуждали ее для списка типов, но ясно, как ее организовать для списка элементов.}
Работает он так: если нам был дан список элементов
\begin{cppcode}[\nonumbers]
FList = AL::List<f0, f2, f2>;
\end{cppcode}
то после применения операции \verb"Enum", мы получим список типов:
\begin{cppcode}
TL::List<Pair<0, f0>, Pair<1, f1>, Pair<2, f2>>;
\end{cppcode}

\item Теперь когда у нас есть пары нам бы хотелось подставить воспользоваться pack-ами внутри \verb"FList".
Но в функции \verb"call" мы не видим его содержимое.
Чтоб увидеть содержимое, надо воспользоваться вспомогательным шаблоном, где содержимое можно просмотреть по явной специализации.
Именно для этого мы передаем вызов следующему шаблону, а точнее вызываем у него метод \verb"exec" и в него форвардим аргументы с помощью perfect forwarding (см.~\ref{section::InitClass} пункт~\ref{section::InitClass::PerfForward}).
\end{enumerate}
Теперь вспомогательный шаблон
\begin{cppcode}

template<class FList>
struct CascadeUnpack {
  static_assert(Logic::False<FList>, "Internal error of cascade_call!");
};

template<template<class...> class TList, auto... fs, int... inds>
struct CascadeUnpack<TList<EL::Pair<inds, fs>...>> {
  template<class... Args>
  constexpr static int exec(Args&&... args) {
    int counter = -1;
    (void)(... || WF<fs, inds>::exec(counter, std::forward<Args>(args)...));
    return counter;
  }
};
\end{cppcode}
Давайте поясню что происходит.
\begin{enumerate}
\item Общий шаблон принимает что угодно, кроме \verb"Enum" (потому что вторая специализация принимает в точности \verb"Enum").
А значит в ней я сообщаю об ошибках.
Причем это ошибка именно \verb"cascade_call", потому что я не хочу, чтобы пользователь этой функции видел имена вспомогательных классов.

\item В специализации \verb"args" -- это переданные аргументы.
Мы заводим счетчик \verb"counter" с дефолтным значением \verb"-1".
И потом уже происходит магия вызова функций через fold-expression.

\item Обратите внимание, что внутри fold-expression мы вызываем не исходные функции \verb"fs", а специальный wrapper
\begin{cppcode}
WF<fs, inds>::exec( ... )
\end{cppcode}
Тут \verb"WF<fs, inds>" -- это шаблонный тип, а \verb"exec" -- статическая функция внутри него.
Эта функция принимает все аргументы \verb"args" и дополнительно \verb"counter" по ссылке и внутри пишет значение из \verb"inds", если \verb"fs" вернула \verb"true".
Обратите внимание, что функции и индексы передаются через шаблонные параметры.
И нам придется обсудить еще один wrapper.

\item Само fold-expression выглядит так:
\begin{cppcode}[\nonumbers]
(... || WF<fs, inds>::exec(counter, std::forward<Args>(args)...));
\end{cppcode}
Что грубо говоря превращается в:
\begin{cppcode}[\nonumbers]
(WF<f0, 0>::exec(counter, std::forward<Args>(args)...) ||
 WF<f1, 1>::exec(counter, std::forward<Args>(args)...) ||
 WF<f2, 2>::exec(counter, std::forward<Args>(args)...));
\end{cppcode}
Обратите внимание, что это выражение возвращает значение типа \verb"bool", но мы его не используем.
И чтобы подавить предупреждения компилятора, я использую \verb"(void)" перед возвращаемым значением.

\item Ну и после вычисления fold-expression возвращаем сохраненное значение \verb"counter".
\end{enumerate}
Ну и теперь последний рывок -- внутренний wrapper.
\begin{cppcode}
template<auto f, int ind>
struct FW {
  template<class... Args>
  constexpr static bool exec(int& counter, Args&&... args) {
    if (f(std::forward<Args>(args)...)) {
      counter = ind;
      return true;
    }
    return false;
  }
};
\end{cppcode}
Тут вроде без сюрпризов, мы запускаем функцию \verb"f" на ее наборе аргументов внутри \verb"if".
Если получили \verb"true", то прежде чем выйти из функции, мы записываем в \verb"counter" текущее значение индекса \verb"ind" и уже потом выходим.
А если \verb"false", то так \verb"false" и возвращаем и не паримся.

\subsubsection{Добавляем диспетчеризацию}

\paragraph{Что хотим}

Теперь давайте усложним задачу.
Пусть у нас есть список переменных заданных заранее как и прежде:
\begin{cppcode}[\nonumbers]
A a; B b; C c; D d;
\end{cppcode}
Но теперь каждая из функций принимает не обязательно все эти переменные в качестве аргументов.
И при этом не обязательно в том же порядке.
Но гарантируется, что \verb"A", \verb"B", \verb"C", \verb"D" -- разные типы и каждая функция принимает каждую переменную не более одного раза.
Правда функция может принимать переменную не только по значению, но и по ссылке и константной ссылке, но НЕ rvalue ссылке (это тоже можно было бы обойти, но я не хочу заморачиваться).
\begin{cppcode}[\nonumbers]
bool f0(A, const C&);
bool f1(C, B);
bool f2(A&, D, const C&);
\end{cppcode}
И мы хотим использовать тот же самый синтаксис что и выше:
\begin{cppcode}[\nonumbers]
using FList = AL::List<f0, f1, f2>;
int res = cascade_auto_call<FList>(a, b, c, d);
\end{cppcode}
И последняя строчка вызовет функции слева направо, подставляя аргументы куда надо.
Для имплементации такой версии нам надо добавить автоматическую диспетчеризацию для аргументов.
Ее мы как раз обсуждали в разделе~\ref{section::ArgDispatcher}.

\paragraph{Имплементация}

Если посмотреть внимательно, то мы увидим, что единственное, чем отличаются две функции по каскадному вызову -- это версией wrapper-а.
Давайте я покажу, как будет выглядеть этот wrapper для автоматической диспетчиризации.
\begin{cppcode}
template<auto f, int ind>
struct FWAuto {
  template<class... Args>
  constexpr static bool exec(int& counter, Args&&... args) {
    if (auto_call(f, std::forward<Args>(args)...)) {
      counter = ind;
      return true;
    }
    return false;
  }
};
\end{cppcode}
То есть теперь нам надо обернуть вызов функции внутрь \verb"auto_call", чтоб она расставила аргументы как нам надо.
А чтобы переиспользовать код, давайте просто сделаем wrapper еще одним параметром для двух вспомогательных классов.
А именно
\begin{cppcode}
template<class FList, template<auto, int> class WF>
struct CascadeUnpack {
  static_assert(Logic::False<FList>, "Internal error of cascade_call!");
};

template<template<class...> class TList, auto... fs, int... inds,
         template<auto, int> class WF>
struct CascadeUnpack<TList<EL::Pair<inds, fs>...>, WF> {
  template<class... Args>
  constexpr static int exec(Args&&... args) {
    int counter = -1;
    (void)(... || WF<fs, inds>::exec(counter, std::forward<Args>(args)...));
    return counter;
  }
};
\end{cppcode}
Тут мы принимаем \verb"WF" как шаблонный параметр, а не используем конкретный шаблон.
Чтобы мы могли передать сюда в распаковку этот шаблон, про него должен быть в курсе и \verb"CascadeCall" класс.
В него мы тоже добавим этот параметр
\begin{cppcode}
template<class FList, template<auto, int> class WF,
         bool isCorrect = AL::isList<FList>>
struct CascadeImpl {
  static_assert(
      AL::isList<FList>,
      "The template argument of cascade_call MUST be an element list of "
      "functions returning bool!");
};

template<class FList, template<auto, int> class WF>
struct CascadeImpl<FList, WF, true> {
  template<class... Args>
  constexpr static int call(Args&&... args) {
    return CascadeUnpack<AL::Enum<FList>, WF>::exec(
        std::forward<Args>(args)...);
  }
};
\end{cppcode}
И мы просто передаем его внутрь.
А теперь уже наши две версии каскадного вызова будут выглядеть так:
\begin{cppcode}
template<class FList, class... Args>
constexpr int cascade_call(Args&&... args) {
  return detail::CascadeImpl<FList, detail::FWrap>::call(
      std::forward<Args>(args)...);
}

template<class FList, class... Args>
constexpr int cascade_auto_call(Args&&... args) {
  return detail::CascadeImpl<FList, detail::FWrapAuto>::call(
      std::forward<Args>(args)...);
}
\end{cppcode}
Вот и все.
Теперь можно последовательно вызывать функции друг за другом до нужного момента и не заботиться о том, где какие аргументы стоят.
